%% ma: L12-B08-0169
%% lop: 12
%% chuong: 2
%% bai: 8
%% dang: 12.8.53
%% loai: TN4
%% muc: VD
%% dap_an: "A"
%% nguon: "CĐ2.4.pdf giáo viên gửi 27/10/2026 (M07, Lê Minh Khôi – Nguyễn Anh Khôi, Chương II Chủ đề 4 Các bài toán về tâm tỉ cự), Câu 19"
%% kien_thuc: "Bài 8 (khoảng cách hai điểm, biểu thức tọa độ), Bài 6 (bình phương vô hướng, tâm tỉ cự)"
%% trang_thai: chinh-thuc
%% ghi_chu: "Giáo viên đã duyệt 27/10/2026"
\begin{ex}
Trong không gian $Oxyz$, cho ba điểm $A(1;0;1)$, $B(0;1;-2)$, $C(-3;2;0)$. Tọa độ của điểm $L$ thỏa mãn $3LA^2-2LB^2+LC^2$ đạt giá trị nhỏ nhất là
\choice{\True $L\left(0;0;\dfrac72\right)$}{$L\left(0;0;-\dfrac72\right)$}{$L\left(10;0;\dfrac72\right)$}{$L\left(0;1;\dfrac72\right)$}
\loigiai{Gọi $I$ thỏa $3\overrightarrow{IA}-2\overrightarrow{IB}+\overrightarrow{IC}=\vec 0\Rightarrow I=\dfrac{3A-2B+C}2=\left(0;0;\dfrac72\right)$.

Biểu thức bằng $2LI^2+$ hằng số, nhỏ nhất khi $L\equiv I$.}
\end{ex}
